\chapter{Conclusions}
\label{ch:conclusion}

This thesis investigated the implications of anycast routing for data sovereignty, international data transfers, and geographic control over Internet communications. The work combined the design of new measurement methodologies, large-scale empirical analyses of operational services, and the development of a mitigation mechanism aimed at reducing jurisdictional exposure in anycast deployments.

The results demonstrate that anycast infrastructures introduce important challenges for transparency and geographic predictability, particularly in regulatory environments where the jurisdiction of personal data destination is a relevant consideration. At the same time, the thesis shows that these challenges can be measured, analysed, and partially mitigated through appropriate architectural mechanisms.

This chapter summarises the main contributions of the thesis, evaluates the fulfilment of the research objectives, discusses the principal limitations of the conducted work, and outlines several directions for future research.

\section{Research Objectives and Main Findings}
\label{sc:conclusion_objectives}

Anycast routing plays a fundamental role in modern Internet infrastructures, yet its implications for data sovereignty, international data transfers, and geographic control have received comparatively limited attention. The research presented in this thesis addresses this gap by studying the geographic behaviour of anycast communications and its consequences for transparency, regulatory compliance, and jurisdictional exposure.

To address this gap, the first objective of the thesis was to determine whether the actual destination country of communications towards anycast infrastructures can be inferred with sufficient accuracy. The results presented in Chapter~\ref{ch:hunter} demonstrate that this objective was successfully achieved through the design and validation of Hunter, a methodology capable of estimating the geographic destination of anycast communications using active network measurements and routing analysis techniques.

The second objective was to assess the extent to which anycast infrastructures participate in international personal data transfers and whether such transfers are transparently disclosed to users. The large-scale analysis presented in Chapter~\ref{ch:anycast_experiments} revealed that communications involving personal data are frequently directed towards anycast infrastructures and that these communications often reach jurisdictions different from those expected by users and application developers. The results further showed that many of these destinations are not explicitly disclosed in privacy policies, creating transparency challenges in environments subject to geographic data processing requirements.

A third objective was to understand whether the observed jurisdictional exposure is an isolated phenomenon affecting specific applications or a broader characteristic of operational anycast deployments. The global measurement campaigns conducted throughout this research demonstrated that cross-border routing and jurisdictional variability are common properties of anycast systems, affecting a large fraction of the analysed infrastructures and geographic regions.

Finally, the thesis sought to explore whether geographic control could be improved without abandoning the operational benefits of anycast. The RCA-DNS architecture presented in Chapter~\ref{ch:rcadns} showed that introducing geographic constraints during service selection can significantly reduce jurisdictional exposure while preserving the scalability and resilience properties that motivate the deployment of anycast infrastructures.

The results obtained throughout this research demonstrate that geographic destination should not be regarded as a static property of anycast communications. Instead, it emerges as a dynamic characteristic influenced by routing decisions, network conditions, and infrastructure deployment strategies. Consequently, any meaningful assessment of data sovereignty and international transfers in anycast environments requires both visibility into routing behaviour and mechanisms capable of influencing geographic service selection.

\section{Scientific Contributions}
\label{sc:conclusion_contributions}

The research presented in this thesis contributes to the study of anycast infrastructures, geographic control, and data sovereignty through three complementary lines of work spanning measurement, empirical analysis, and system design.

The first contribution is Hunter, a methodology for inferring the geographic destinations reached by communications towards anycast infrastructures. Hunter combines traceroute analysis, last-hop identification, latency-based multilateration, and geographic inference techniques to estimate the destination country associated with operational anycast deployments. The validation results demonstrate that the proposed methodology achieves high levels of accuracy while remaining applicable at Internet scale. By providing visibility into the geographic behaviour of anycast communications, Hunter enables analyses that were previously difficult to perform systematically.

Building upon this capability, the second contribution consists of a large-scale empirical analysis of international personal data transfers involving anycast infrastructures in Android mobile ecosystems. By combining dynamic traffic analysis, privacy policy inspection, and anycast destination inference, the study revealed the prevalence of international transfers through anycast services and quantified the discrepancy between the jurisdictions reached in practice and those disclosed to users. The analysis further demonstrated the significant role played by Third-Party Libraries in generating these transfers and showed that jurisdictional exposure is not restricted to specific applications or providers, but rather emerges as a systemic characteristic of operational anycast deployments.

The third contribution is the design and experimental evaluation of Region-Constrained Anycast over DNS (RCA-DNS), a practical architecture for introducing geographic constraints into anycast service deployments. The proposed approach combines DNS-based service selection with geographically scoped anycast infrastructures, enabling service providers to influence the jurisdictions reached by user communications while preserving the operational benefits of anycast. The experimental evaluation showed that geographic confinement can be achieved while maintaining latency levels compatible with modern Internet services, demonstrating the practical feasibility of jurisdiction-aware anycast deployments.

Beyond these specific contributions, the thesis establishes a broader connection between Internet routing, geographic control, and regulatory compliance. The results demonstrate that routing decisions can directly influence the jurisdictions where communications are processed and that geographic destination should therefore be considered an operational property of distributed systems rather than a static characteristic of service deployments. This perspective extends the traditional understanding of anycast and highlights the importance of incorporating geographic considerations into the design, deployment, and evaluation of future Internet infrastructures.

\section{Limitations}
\label{sc:conclusion_limitations}

The destination inference methodology developed in this thesis relies on active network measurements and geolocation information. As discussed throughout Chapter~\ref{ch:hunter}, although Hunter achieved high accuracy and precision during its validation, destination inference remains inherently dependent on the quality of the underlying measurement data and geolocation sources. Internet routing dynamics, incomplete measurement visibility, and inaccuracies in public geolocation databases may affect individual inferences. Nevertheless, the large-scale nature of the experiments and the repeated measurements performed throughout the study mitigate the impact of isolated errors on the overall conclusions.

The large-scale analysis of Android applications was based on dynamic execution techniques. As with any dynamic analysis methodology, complete behavioural coverage cannot be guaranteed. Certain application functionalities may only be triggered under specific user interactions, authentication requirements, geographic conditions, or runtime contexts that were not reproduced during the experiments. Consequently, the observed communications should be interpreted as a lower bound of the actual data transfers performed by the analysed applications and Third-Party Libraries.

The privacy policy analysis focused on identifying the jurisdictions disclosed as possible destinations for international transfers. While the methodology enabled a systematic comparison between declared and observed destinations, privacy policies often contain ambiguous language, generic statements, or references to external providers that complicate automated interpretation. Furthermore, privacy documentation may evolve over time, potentially introducing changes after the measurements were conducted.

The global measurements performed using Hunter were necessarily constrained by the availability and geographic distribution of measurement vantage points. Although RIPE Atlas provided broad coverage across the studied regions, certain countries and networks remain under-represented. Consequently, the observed routing behaviours may not capture every possible destination reachable through a given anycast deployment.

The evaluation of RCA-DNS focused on demonstrating the feasibility of introducing geographic constraints into anycast service deployments. The experimental validation was conducted using a limited number of service replicas deployed within Europe and therefore does not fully represent the complexity of large-scale commercial deployments. Moreover, the evaluation concentrated on latency and geographic confinement, leaving other operational aspects such as large-scale scalability, resilience under failures, and long-term behaviour outside the scope of the study. Additional experiments involving larger infrastructures, multiple geographic regions, and longer observation periods would be necessary to assess the behaviour of the proposed architecture under highly dynamic operational conditions.

Despite these limitations, the methodologies, measurements, and analyses presented throughout this thesis consistently reveal a previously underexplored relationship between anycast routing, international data transfers, and geographic control. Although the exact magnitude of the observed phenomena may vary under different experimental conditions, the results obtained across multiple datasets, measurement campaigns, and methodological approaches support the robustness of the main conclusions. Taken together, these findings provide strong evidence that geographic destination is a dynamic property of anycast communications and that routing behaviour must be considered when analysing data sovereignty and jurisdictional exposure in modern Internet infrastructures.

\section{Future Research Directions}
\label{sc:conclusion_future}

The destination inference methodology developed in this thesis can be extended in several directions. Although Hunter demonstrated high accuracy in identifying the destination countries of anycast communications, additional improvements could be achieved by incorporating new measurement sources, routing datasets, or geolocation techniques. Future work could also investigate mechanisms for continuously monitoring anycast deployments and automatically detecting changes in their geographic footprint over time. Such capabilities would be particularly valuable for organisations seeking to assess jurisdictional exposure in dynamic network environments.

The large-scale analysis presented in this thesis focused on the Android ecosystem due to its widespread adoption and extensive reliance on cloud-based services and Third-Party Libraries. However, the same methodology could be applied to other environments, including web applications, Internet of Things platforms, content delivery systems, and cloud-native services. Extending the analysis to these domains would provide a broader understanding of how anycast infrastructures influence cross-border communications across the Internet.

The transparency challenges identified throughout Chapter~\ref{ch:anycast_experiments} also motivate further research into mechanisms capable of communicating geographic processing locations and jurisdictional exposure more effectively. Current privacy policies typically describe international transfers through static and high-level statements that are poorly suited to the dynamic nature of anycast routing. New approaches capable of representing destination variability and routing-induced geographic changes could improve the information available to users and organisations.

The relationship between anycast infrastructures and emerging regulatory requirements related to data sovereignty constitutes another promising area for future work. The results obtained throughout this thesis suggest that geographic destination should be considered a dynamic operational property rather than a fixed characteristic of a service. Consequently, compliance assessment methodologies, auditing processes, and regulatory frameworks may need to evolve in order to account for routing-induced geographic variability.

The RCA-DNS architecture presented in Chapter~\ref{ch:rcadns} represents an initial step towards jurisdiction-aware anycast deployments. Evaluating the approach under larger-scale operational conditions involving additional geographic regions, larger numbers of replicas, and longer observation periods would provide a more comprehensive understanding of the trade-offs between geographic confinement, performance, availability, and operational complexity.

The policy model implemented by RCA-DNS could also be extended beyond the regional constraints explored in this thesis. Future systems may support dynamic geographic policies that adapt to changes in service requirements, infrastructure availability, contractual obligations, or regulatory conditions. Such extensions could allow service providers to express geographic requirements at different levels of granularity, ranging from broad regions to individual countries or specific legal jurisdictions.

Alternative approaches to geographic control also deserve further investigation. Integrating geographic awareness directly into routing decisions, traffic engineering policies, cloud orchestration platforms, or service deployment strategies may provide additional levels of control over the jurisdictions reached by Internet communications. Exploring these approaches could contribute to the development of Internet infrastructures capable of simultaneously satisfying performance, resilience, and geographic governance requirements.

More broadly, the results of this thesis suggest that geographic awareness is likely to become an increasingly important design consideration for distributed Internet services. As global infrastructures continue to expand and regulatory requirements surrounding data sovereignty evolve, understanding and controlling the geographic behaviour of networked systems will remain an important and challenging area of research.

\section{Closing Remarks}
\label{sc:conclusion_closing}

Anycast has become a fundamental building block of modern Internet infrastructures. Its ability to improve performance, scalability, and resilience has led to its widespread adoption by cloud providers, content delivery networks, security services, and countless online platforms. However, as this thesis has shown, the geographic consequences of anycast routing extend beyond traditional networking concerns and directly influence where communications are processed and, ultimately, under which jurisdictions data may be exposed.

The development of Hunter demonstrated that the geographic destinations reached by anycast communications can be systematically inferred and analysed. The large-scale measurements conducted throughout this thesis revealed that cross-border routing and jurisdictional variability are common characteristics of operational anycast deployments, affecting both mobile applications and the third-party ecosystems on which they depend. The results further showed that these geographic dynamics frequently remain invisible to users and are often insufficiently reflected in existing transparency mechanisms.

These findings suggest that geographic location should not be regarded as a static attribute of Internet services. Instead, it emerges from the interaction between infrastructure deployment, routing decisions, and service architectures. As regulatory requirements related to data sovereignty and international data transfers continue to evolve, understanding these interactions will become increasingly important for service providers, regulators, and researchers alike.

The RCA-DNS architecture presented in this thesis illustrates that the challenges introduced by anycast are not necessarily incompatible with geographic control. By introducing geographic awareness into the service selection process, it is possible to reduce jurisdictional exposure while preserving many of the operational properties that motivate the use of anycast infrastructures.

More broadly, this thesis highlights that Internet routing decisions can have direct implications for transparency, data sovereignty, and regulatory compliance. Geographic destination is not simply a property of where infrastructure is deployed but also of how communications are routed through that infrastructure. Recognising this relationship is essential for understanding the real geographic behaviour of distributed services and for designing future Internet infrastructures capable of balancing performance, resilience, and geographic governance requirements.
